\documentclass[11pt,fontset=fandol]{ctexart}
\usepackage[a4paper,margin=23mm]{geometry}
\usepackage{amsmath,amssymb,mathtools,bm,mathrsfs,esint}
\usepackage{graphicx,longtable,booktabs,array,calc}
\usepackage[normalem]{ulem}
\usepackage{xcolor,hyperref}
\hypersetup{hidelinks}
\setcounter{secnumdepth}{-1}
\newcounter{none}
\providecommand{\tightlist}{\setlength{\itemsep}{0pt}\setlength{\parskip}{0pt}}
\providecommand{\passthrough}[1]{#1}
\usepackage{amssymb,mathrsfs,esint,bm,newunicodechar}
\allowdisplaybreaks
\setlength{\emergencystretch}{3em}
\providecommand{\llangle}{\langle\!\langle}
\providecommand{\rrangle}{\rangle\!\rangle}
\newunicodechar{∞}{\ensuremath{\infty}}
\newunicodechar{∈}{\ensuremath{\in}}
\newunicodechar{∑}{\ensuremath{\sum}}
\newunicodechar{∫}{\ensuremath{\int}}
\newunicodechar{≤}{\ensuremath{\leq}}
\newunicodechar{≥}{\ensuremath{\geq}}
\newunicodechar{≠}{\ensuremath{\neq}}
\newunicodechar{→}{\ensuremath{\to}}
\newunicodechar{⇒}{\ensuremath{\Rightarrow}}
\newunicodechar{α}{\ensuremath{\alpha}}
\newunicodechar{β}{\ensuremath{\beta}}
\newunicodechar{γ}{\ensuremath{\gamma}}
\newunicodechar{λ}{\ensuremath{\lambda}}
\newunicodechar{φ}{\ensuremath{\varphi}}
\newunicodechar{ε}{\ensuremath{\varepsilon}}
\newunicodechar{σ}{\ensuremath{\sigma}}
\newunicodechar{θ}{\ensuremath{\theta}}
\newunicodechar{δ}{\ensuremath{\delta}}
\newunicodechar{η}{\ensuremath{\eta}}
\newunicodechar{π}{\ensuremath{\pi}}
\newunicodechar{ξ}{\ensuremath{\xi}}
\newunicodechar{Σ}{\ensuremath{\Sigma}}
\newunicodechar{Δ}{\ensuremath{\Delta}}
\newunicodechar{∂}{\ensuremath{\partial}}
\newunicodechar{∀}{\ensuremath{\forall}}
\newunicodechar{∃}{\ensuremath{\exists}}
\newunicodechar{√}{\ensuremath{\surd}}
\newunicodechar{−}{\ensuremath{-}}
\newunicodechar{μ}{\ensuremath{\mu}}
\newunicodechar{ρ}{\ensuremath{\rho}}
\newunicodechar{τ}{\ensuremath{\tau}}
\newunicodechar{ω}{\ensuremath{\omega}}
\newunicodechar{Ω}{\ensuremath{\Omega}}
\newunicodechar{χ}{\ensuremath{\chi}}
\newunicodechar{ψ}{\ensuremath{\psi}}
\newunicodechar{Γ}{\ensuremath{\Gamma}}
\newunicodechar{Φ}{\ensuremath{\Phi}}
\newunicodechar{Λ}{\ensuremath{\Lambda}}
\newunicodechar{ℝ}{\ensuremath{\mathbb{R}}}
\newunicodechar{ℤ}{\ensuremath{\mathbb{Z}}}
\newunicodechar{ℕ}{\ensuremath{\mathbb{N}}}
\newunicodechar{△}{\ensuremath{\triangle}}
\newunicodechar{⊗}{\ensuremath{\otimes}}
\begin{document}
\section{2017年杭州电子科技大学601数学分析考研真题参考答案}\label{ux5e74ux676dux5ddeux7535ux5b50ux79d1ux6280ux5927ux5b66601ux6570ux5b66ux5206ux6790ux8003ux7814ux771fux9898ux53c2ux8003ux7b54ux6848}

来源：\texttt{HDU/601数学分析/2017/2017年杭州电子科技大学601数学分析考研真题参考答案.pdf}。原
PDF 页序：1, 2, 3, 4, 5, 6, 7, 8。

\section{杭州电子科技大学2017年数学分析真题解析}\label{ux676dux5ddeux7535ux5b50ux79d1ux6280ux5927ux5b662017ux5e74ux6570ux5b66ux5206ux6790ux771fux9898ux89e3ux6790}

\begin{quote}
助手整理依据：保留全部题号和原解答思路，依据打印题面及同题前后等式恢复被水印破坏的公式。一5原答案算的是
\(\partial z/\partial x\)，本稿保留该式并另给出打印题目实际要求的
\(\partial z/\partial y\)。一7系数的数字 \(9\) 被OCR误识别为字母
\(g\)，按打印题面恢复。四1的负二阶导数按实际求导恢复为
\(f''\)，不沿用OCR误写的 \(f'\)。
\end{quote}

\subsection{一、填空题（8题，每题5分，共40分）}\label{ux4e00ux586bux7a7aux98988ux9898ux6bcfux98985ux5206ux517140ux5206}

\subsubsection{1. 极限}\label{ux6781ux9650}

\[
\lim_{x\to0}\frac{\ln(1+x^2)}{x\tan2x}
=\lim_{x\to0}\frac{x^2}{2x^2}=\frac12.
\]

\subsubsection{2. 间断点}\label{ux95f4ux65adux70b9}

对题给分段表达式，可能的异常点为 \(0,1\)。单侧极限为 \[
\lim_{x\to1^+}e^{1/(x-1)}=+\infty,\qquad
\lim_{x\to1^-}e^{1/(x-1)}=0,
\] 以及 \[
\lim_{x\to0^-}x\sin(1/x)=0,\qquad
\lim_{x\to0^+}e^{1/(x-1)}=e^{-1}.
\] 所以跳跃点为 \(x=0\)；\(x=1\) 是无穷间断点。严格说原分段式未有效定义
\(f(0)\) 和
\(f(1)\)，可将它们理解为表达式的间断点位置；任意补定义点值均不改变这两种单侧极限的分类。将第二支在零点连续延拓为
\(0\) 是常用约定，未将该补充约定冒充原题。

\subsubsection{3. 奇偶性积分}\label{ux5947ux5076ux6027ux79efux5206}

\(\ln(1+x^4)\sin^7x\) 为奇函数，\(|x|e^{|x|}\) 为偶函数，故 \[
\begin{aligned}
\int_{-1}^1[\ln(1+x^4)\sin^7x+|x|e^{|x|}]\,dx
&=2\int_0^1xe^x\,dx\\
&=2[(x-1)e^x]_0^1=2.
\end{aligned}
\]

\subsubsection{4.
积分方程与高阶导数}\label{ux79efux5206ux65b9ux7a0bux4e0eux9ad8ux9636ux5bfcux6570}

由 \(\int_0^tf(u)\,du=f(t)-1\)，首先得到 \(f\)
连续，再由微积分基本定理得到 \(f'=f\)。故 \(f(t)=Ce^t\)。令 \(t=0\) 得
\(f(0)=1\)，所以 \(C=1\)，\(f(t)=e^t\)，进而 \(f^{(n)}(t)=e^t\)。

\subsubsection{5.
隐函数偏导数}\label{ux9690ux51fdux6570ux504fux5bfcux6570}

令 \(F(x,y,z)=e^{-xy}+2z-e^z\)，在 \(F_z=2-e^z\ne0\) 的点，有 \[
\frac{\partial z}{\partial y}=-\frac{F_y}{F_z}
=\frac{xe^{-xy}}{2-e^z}.
\] 这是打印题面所求。原参考答案实际计算的另一个偏导数为 \[
\frac{\partial z}{\partial x}=-\frac{F_x}{F_z}
=-\frac{-ye^{-xy}}{2-e^z}=\frac{ye^{-xy}}{2-e^z}.
\] 两式均成立，但不能把第二式当作对第一式的回答。

\subsubsection{6. 积分换序}\label{ux79efux5206ux6362ux5e8f}

原区域为 \(0\le x\le1\)、\(x^2\le y\le2-x\)。水平截面在 \(y=1\)
处分段，因此 \[
\int_0^1dx\int_{x^2}^{2-x}f(x,y)\,dy
=\int_0^1dy\int_0^{\sqrt y}f(x,y)\,dx
+\int_1^2dy\int_0^{2-y}f(x,y)\,dx.
\]

\subsubsection{7.
幂级数收敛半径}\label{ux5e42ux7ea7ux6570ux6536ux655bux534aux5f84}

设 \(a_n=[9^n+(-4)^n]/n\)。由 \[
R=\lim_{n\to\infty}\sqrt{\frac{a_n}{a_{n+1}}}
=\lim_{n\to\infty}\sqrt{\frac{n+1}{n}\frac{9^n[1+(-4/9)^n]}{9^{n+1}[1+(-4/9)^{n+1}]}}
=\frac13.
\] 根号来自指数为 \(2n\)，不能遗漏。

\subsubsection{8. 全微分}\label{ux5168ux5faeux5206}

对 \(y\ne0\)， \[
dz=\frac{d(x/y)}{1+(x/y)^2}
=\frac{(y\,dx-x\,dy)/y^2}{1+x^2/y^2}
=\frac{y\,dx-x\,dy}{x^2+y^2}.
\]

\subsection{二、计算题（4题，共30分）}\label{ux4e8cux8ba1ux7b97ux98984ux9898ux517130ux5206}

\subsubsection{1. 极限}\label{ux6781ux9650-1}

\[
\begin{aligned}
\lim_{x\to0}\frac{2^x+3^x-2}{\arcsin x}
&=\lim_{x\to0}\frac{2^x+3^x-2}{x}\\
&=\lim_{x\to0}(2^x\ln2+3^x\ln3)=\ln6.
\end{aligned}
\]

\subsubsection{2. 分部积分}\label{ux5206ux90e8ux79efux5206}

\[
\begin{aligned}
\int x^2\arctan x\,dx
&=\frac13\int\arctan x\,d(x^3)\\
&=\frac{x^3}{3}\arctan x-\frac13\int\frac{x^3}{1+x^2}\,dx\\
&=\frac{x^3}{3}\arctan x-\frac13\int\left(x-\frac{x}{1+x^2}\right)dx\\
&=\frac{x^3}{3}\arctan x-\frac{x^2}{6}+\frac16\ln(1+x^2)+C.
\end{aligned}
\] 原OCR把分部积分首项的系数 \(1/3\) 识作
\(\pi/3\)，这里按实际求导恢复。

\subsubsection{3. 反常积分}\label{ux53cdux5e38ux79efux5206}

令 \(x=\sin t\)，得 \[
\int_0^1\frac{x^5}{\sqrt{1-x^2}}\,dx
=\int_0^{\pi/2}\sin^5t\,dt=\frac{4\cdot2}{5\cdot3}=\frac8{15}.
\]

\subsubsection{4.
上半球面的通量}\label{ux4e0aux534aux7403ux9762ux7684ux901aux91cf}

按半球半径 \(a>0\) 理解。添加底面 \(\Sigma_1=\{z=0,x^2+y^2\le a^2\}\)
并取下侧，则底面通量为零。由 Gauss 公式，设
\(\Omega=\{x^2+y^2+z^2\le a^2,z\ge0\}\)， \[
\begin{aligned}
\iint_\Sigma x^3\,dy\,dz+y^3\,dz\,dx+z^3\,dx\,dy
&=\iint_{\Sigma\cup\Sigma_1}x^3\,dy\,dz+y^3\,dz\,dx+z^3\,dx\,dy\\
&=\iiint_\Omega\left(\frac{\partial x^3}{\partial x}+\frac{\partial y^3}{\partial y}+\frac{\partial z^3}{\partial z}\right)dV\\
&=3\iiint_\Omega(x^2+y^2+z^2)\,dV\\
&=\frac32\iiint_{x^2+y^2+z^2\le a^2}(x^2+y^2+z^2)\,dV\\
&=\frac32\cdot4\pi\int_0^ar^4\,dr=\frac65\pi a^5.
\end{aligned}
\]
原OCR中的多重积分号、区域下标和微分符号已按散度公式恢复。若原题允许任意实数
\(a\)，几何半径为 \(|a|\)，结果相应为 \(6\pi|a|^5/5\)。

\subsection{三、综合题（4题，每题10分，共40分）}\label{ux4e09ux7efcux5408ux98984ux9898ux6bcfux989810ux5206ux517140ux5206}

\subsubsection{1. 一致收敛}\label{ux4e00ux81f4ux6536ux655b}

利用 \(|\arctan x|\le\pi/2\) 和
\(1+n^3x^2\ge2n^{3/2}|x|\)，可将原答案受水印损坏的界恢复为 \[
\left|\frac{(-1)^nx\arctan x}{1+n^3x^2}\right|
\le\frac{\pi|x|}{2(1+n^3x^2)}\le\frac\pi{4n^{3/2}}\le\frac\pi{2n^{3/2}}.
\] 在 \(x=0\) 时不等式同样成立。由于 \(\sum\pi/(2n^{3/2})\)
收敛，原级数在 \(\mathbb R\) 上一致且绝对收敛。

\subsubsection{2.
参数幂级数收敛域}\label{ux53c2ux6570ux5e42ux7ea7ux6570ux6536ux655bux57df}

系数比值极限为 \[
\lim_{n\to\infty}\frac{1/(n+p)^p}{1/(n+1+p)^p}=1,
\] 故半径为 \(1\)。\(x=-1\) 时由 Leibniz 判别法收敛；\(x=1\) 时， \[
\lim_{n\to\infty}\frac{1/(n+p)^p}{1/n^p}=1,
\] 所以当且仅当 \(p>1\) 时收敛。答案为 \(0<p\le1\) 时
\([-1,1)\)，\(p>1\) 时 \([-1,1]\)。

\subsubsection{\texorpdfstring{3. 求常数
\(k\)}{3. 求常数 k}}\label{ux6c42ux5e38ux6570-k}

有 \[
\lim_{x\to\infty}\frac{x+k}{x-k}=1,\qquad
\lim_{x\to\infty}x\left(\frac{x+k}{x-k}-1\right)
=\lim_{x\to\infty}\frac{2kx}{x-k}=2k,
\] 所以左端极限为 \(e^{2k}\)。由中值定理，存在 \(\xi_x\in(x-1,x+1)\)，使
\(f(x+1)-f(x-1)=2f'(\xi_x)\)，且 \(\xi_x\to\infty\)，所以右端极限为
\(2e\)。于是 \[
e^{2k}=2e,\qquad k=\frac{1+\ln2}{2}.
\]

\subsubsection{4.
约束极值与不等式}\label{ux7ea6ux675fux6781ux503cux4e0eux4e0dux7b49ux5f0f}

在闭三角域 \(x,y,z\ge0,x+y+z=a\) 上，\(f=xy^2z^3\)
连续、边界为零、内部为正，最大值必在内部取得。对 \(g=\ln f\)
使用乘子法，令 \[
L=\ln(xy^2z^3)+\lambda(x+y+z-a),
\] 驻点方程为 \[
\frac1x+\lambda=0,\quad\frac2y+\lambda=0,\quad
\frac3z+\lambda=0,\quad x+y+z=a.
\] 所以 \(x:y:z=1:2:3\)，唯一候选点为
\((a/6,a/3,a/2)\)，因此它是最大值点， \[
f_{\max}=\frac a6\frac{a^2}{9}\frac{a^3}{8}=\frac{a^6}{432}.
\] 正变量区域无最小值，其下确界为零。令 \(a=x+y+z\)，即得 \[
xy^2z^3\le\frac{(x+y+z)^6}{432}
=108\left(\frac{x+y+z}{6}\right)^6,
\] 等号当且仅当 \(x:y:z=1:2:3\)。

\subsection{四、证明题（4题，每题10分，共40分）}\label{ux56dbux8bc1ux660eux98984ux9898ux6bcfux989810ux5206ux517140ux5206}

\subsubsection{1. 正弦不等式}\label{ux6b63ux5f26ux4e0dux7b49ux5f0f}

令 \(F(x)=\sin(x/2)-x/\pi\)，\(x\in[0,\pi]\)，有 \(F(0)=F(\pi)=0\)， \[
F''(x)=-\frac14\sin(x/2)<0,\qquad 0<x<\pi.
\] 所以 \(F\)
严格凹，图像在端点弦的上方，\(F(x)>0\)，结论成立。原OCR将此处二阶导数写成一阶导数，已按实际求导恢复。

\subsubsection{2.
严格凹函数的不等式}\label{ux4e25ux683cux51f9ux51fdux6570ux7684ux4e0dux7b49ux5f0f}

\(f''<0\) 保证 \(f'\) 严格递减。不妨设 \(x_1\le x_2\)，则 \[
f(x_1+x_2)-f(x_1)-f(x_2)
=[f(x_1+x_2)-f(x_2)]-[f(x_1)-f(0)].
\] 分别对两个差值用中值定理，存在
\(\xi\in(x_2,x_1+x_2)\)、\(\eta\in(0,x_1)\)，使 \[
f(x_1+x_2)-f(x_2)=x_1f'(\xi),\qquad f(x_1)-f(0)=x_1f'(\eta).
\] 因为 \(\xi>\eta\)，有 \(f'(\xi)<f'(\eta)\)。相减得到
\(f(x_1+x_2)<f(x_1)+f(x_2)\)。

\subsubsection{\texorpdfstring{3. \(x^2\)
不一致连续}{3. x\^{}2 不一致连续}}\label{x2-ux4e0dux4e00ux81f4ux8fdeux7eed}

\(f(x)=x^2\) 在 \((0,\infty)\) 连续。取 \(x_n=\sqrt n\)，则 \[
x_{n+1}-x_n=\frac1{\sqrt{n+1}+\sqrt n}\to0,
\qquad f(x_{n+1})-f(x_n)=1.
\] 因此不满足一致连续的序列判据。

\subsubsection{4. Green
公式及对称性}\label{green-ux516cux5f0fux53caux5bf9ux79f0ux6027}

令
\(I_+=\oint_Lxe^{\sin y}\,dy-ye^{-\sin x}\,dx\)，\(I_-=\oint_Lxe^{-\sin y}\,dy-ye^{\sin x}\,dx\)。

（1）由 Green 公式， \[
\begin{aligned}
I_+&=\iint_D\left[\frac{\partial(xe^{\sin y})}{\partial x}-\frac{\partial(-ye^{-\sin x})}{\partial y}\right]dx\,dy
=\iint_D(e^{\sin y}+e^{-\sin x})\,dx\,dy,\\
I_-&=\iint_D\left[\frac{\partial(xe^{-\sin y})}{\partial x}-\frac{\partial(-ye^{\sin x})}{\partial y}\right]dx\,dy
=\iint_D(e^{-\sin y}+e^{\sin x})\,dx\,dy.
\end{aligned}
\] 由于 \(D\) 关于 \(x=y\) 对称， \[
\iint_De^{\sin y}\,dx\,dy=\iint_De^{\sin x}\,dx\,dy,\qquad
\iint_De^{-\sin y}\,dx\,dy=\iint_De^{-\sin x}\,dx\,dy.
\] 因此 \(I_+=I_-\)。

（2）同样由对称性与 \(e^u+e^{-u}\ge2\)， \[
I_+=\iint_D(e^{\sin x}+e^{-\sin x})\,dx\,dy
\ge\iint_D2\,dx\,dy=2S(D)=2\pi^2.
\]
原OCR把有限正方形积分的区域下标识作无穷远，并将末尾不等号混入求和符号，以上按题给区域和前述等式恢复。
\end{document}
